\begin{textsample}{NEG Dim 2 – 2010\_c4 – Score 6.00 – 2010\_c4/enade\_2010\_med\_c0c86d6464.txt}  \label{ex:f2_neg_293}
CAPÍTULO 4 PERCEPÇÃO DA PROVA As análises feitas neste capítulo tratam das percepções dos estudantes, separados por grupos de Ingressantes e Concluintes, sobre a prova aplicada no ENADE/2010 na Área de Medicina.
Estas percepções foram mensuradas por meio de nove questões que avaliaram desde o grau de dificuldade da prova até o tempo gasto para concluí -la.
As percepções sobre a prova foram relacionadas com o desempenho dos alunos e com a Grande Região de funcionamento do curso.
O desempenho dos alunos foi classificado em dois níveis : quarto inferior e quarto superior.
Para tanto, esse desempenho foi colocado em ordem ascendente.
O percentil 25, P25, também conhecido como primeiro quartil, é a nota de desempenho que deixa um quarto ( 25,0 \% ) dos valores observados abaixo e três quartos acima.
A Figura 1 apresenta uma ilustração destes conceitos.
O quarto inferior de desempenho é composto pelas notas abaixo do primeiro quartil.
Já o percentil 75, P75, também conhecido como terceiro quartil, é o valor para o qual há três quartos ( 75,0 \% ) dos dados abaixo e um quarto acima dele.
O quarto superior de desempenho é composto pelas notas iguais ou acima do terceiro quartil.
O percentil 50, P50, também conhecido como mediana, é o valor que divide as notas em dois conjuntos de igual tamanho.
Vale ressaltar que percentis, quartis e medianas são pontos, e não obrigatoriamente pertencentes ao conjunto original de dados.
Ao passo que os quartos são subconjuntos dos dados originais.
A seguir, serão apresentados Gráficos com resultados selecionados, relativos às nove questões avaliadas por grupos de estudantes.
As Tabelas no Anexo II apresentam os valores absolutos e a distribuição percentual das alternativas válidas das nove questões, segundo o mesmo recorte de desempenho dos alunos e Grande Região de funcionamento do curso. 4.1 GRAU DE DIFICULDADE DA PROVA 4.1.1 Componente de Formação Geral Na questão “ Qual o grau de dificuldade desta prova na parte de Formação Geral? ”, 10,6 \% do grupo de Ingressantes e 12,1 \% do grupo de Concluintes optaram pelas alternativas difícil ou muito difícil.
Entretanto, para mais da metade dos estudantes, tanto Ingressantes ( 59,5 \% ) quanto Concluintes ( 64,9 \% ), o Componente de Formação Geral da prova foi considerado com grau de dificuldade médio ( Gráfico 4.1 e Tabelas II.1 e II.10 no Anexo II ).
O percentual de Ingressantes que consideraram a prova difícil ou muito difícil foi mais elevado no quarto inferior ( 14,1 \% ) do que no superior ( 9,0 \% ), o que também foi verificado no grupo dos Concluintes, com percentual de participantes maior no quarto inferior ( 14,5 \% ), do que no quarto superior ( 11,2 \% ).
Semelhantemente ao ocorrido para os Ingressantes como um todo, para os quartos de menor e maior desempenho deste grupo, a alternativa modal para esta pergunta foi médio, com respectivamente 57,9 \% e Já entre os Concluintes, estas percentagens foram, respectivamente, 65,2 \% e 63,2 \%.
Considerando -se as alternativas escolhidas pelos estudantes de cada Grande Região, a região Norte foi aquela na qual houve maior proporção ( 12,0 \% ), no grupo de Ingressantes, que avaliou o Componente de Formação Geral da prova como difícil ou muito difícil.
Em contrapartida, a região Nordeste foi a que apresentou percentual menor : 8,7 \% desse grupo.
Dentre os Concluintes, a maior incidência dos que consideraram a parte de Formação Geral difícil ou muito difícil estava na região Centro-Oeste ( 14,7 \% ), ao passo que, o menor índice foi verificado na região Nordeste, com 8,9 \% dos participantes desse grupo.
A opção pela alternativa de grau médio por Grande Região, tanto no grupo dos Concluintes quanto no grupo dos Ingressantes, localizou -se em maior medida na região Norte, com 68,3 \% e 62,8 \%, respectivamente. 4.1.2 Componente de Conhecimento Específico Ao responderem à questão “ Qual o grau de dificuldade desta prova na parte de Componente Específico? ”, 66,3 \% do grupo de Ingressantes e 13,9 \% do grupo de Concluintes classificaram -na como difícil ou muito difícil.
Este Componente de Conhecimento Específico da prova foi considerado com grau de dificuldade médio por 30,6 \% dos Ingressantes e por 71,8 \% dos Concluintes ( Gráfico 4.2 e Tabelas II.2 e II.11 no Anexo II ).
Considerando -se a percepção sobre o grau de dificuldade do Componente de Conhecimento Específico da prova de acordo com o desempenho dos estudantes, observa -se que 66,9 \% dos Ingressantes que obtiveram desempenho situado no quarto superior classificaram -na como difícil ou muito difícil.
Dos Ingressantes que obtiveram desempenho situado no quarto inferior, 64,0 \% optaram pela mesma alternativa.
Para os Concluintes, o percentual de estudantes que consideraram o Componente de Conhecimento Específico da prova difícil ou muito difícil foi um pouco maior entre os que obtiveram desempenho no quarto inferior ( 17,1 \% ) do que os que obtiveram desempenho no quarto superior ( 10,4 \% ).
Semelhantemente ao ocorrido para os Ingressantes como um todo, para os quartos de menor e maior desempenho deste grupo, a alternativa modal para esta pergunta foi difícil, com respectivamente 45,4 \% e 53,9 \% dos respondentes.
Já entre os Concluintes, a alternativa modal para a Questão 2 foi médio, com 71,2 \% do quarto inferior e 70,5 \% do quarto superior.
Entre as Grandes Regiões brasileiras, a proporção de Ingressantes que avaliaram o Componente de Conhecimento Específico da prova como difícil ou muito difícil variou de 56,9 \% ( Nordeste ) a 71,2 \% ( Centro-Oeste ).
O grupo de Concluintes avaliou esse componente como difícil ou muito difícil com menor frequência, com o mínimo alcançado na região Sudeste, 12,9 \%, e o máximo na região Centro-Oeste, com 19,4 \% dos respondentes.
O percentual de Ingressantes que classificaram o grau de dificuldade como médio, no Componente de Conhecimento Específico da prova, variou de 25,7 \% ( Centro-Oeste ) a 39,6 \% ( Nordeste ).
Dentre os Concluintes, a percepção do mesmo grau de dificuldade aconteceu com maior frequência, variando de 69,6 \% ( Centro-Oeste ) a 74,1 \% ( Nordeste ). 4.2 EXTENSÃO DA PROVA EM RELAÇÃO AO TEMPO TOTAL Indagados quanto à extensão da prova, em relação ao tempo total oferecido para a sua resolução, os estudantes apontaram, com maior incidência, a alternativa que considerava a extensão adequada, para ambos os grupos e todas as agregações avaliadas ( Gráfico 4.3 e Tabelas II.3 e II.12 no Anexo II ).
No grupo de Ingressantes, a maioria, 69,2 \%, respondeu ser a extensão da prova adequada, ao passo que, no grupo de Concluintes, este valor foi menor, 63,1 \%.
Para 27,1 \% dos Ingressantes o exame foi longo ou muito longo.
Essa mesma consideração foi feita por 31,4 \% dos Concluintes.
No grupo de Ingressantes, a maioria, tanto para o quarto de desempenho inferior ( 64,7 \% ) quanto para o quarto de melhor desempenho ( 70,0 \% ) respondeu ser a extensão da prova adequada.
Já entre os Concluintes, esses valores foram um pouco menores, situados em cada quarto em, respectivamente, 60,0 \% e 66,7 \%.
Considerando -se ainda os quartos extremos de desempenho, observa -se que a proporção de estudantes que consideraram a prova longa ou muito longa, em relação ao tempo total destinado à sua resolução, acompanhou a percepção de Ingressantes e Concluintes.
Para o grupo de Ingressantes, entre os estudantes no quarto de desempenho inferior, 30,3 \% declararam ter esta percepção, ao passo que entre aqueles no quarto de melhor desempenho, esta proporção foi de 26,5 \%.
O mesmo pode ser constatado entre os Concluintes : 35,5 \% no quarto de menor desempenho e 27,3 \% no quarto de maior desempenho declararam ter tido esta percepção.
Na análise por Grandes Regiões, o percentual de Ingressantes que responderam ser a extensão da prova adequada variou de 64,7 \% no Nordeste a 76,1 \%, no Centro-Oeste.
Já entre os Concluintes, estes valores estavam compreendidos entre 61,3 \% ( Nordeste ) e 68,5 \% ( Centro-Oeste ).
Aqueles Ingressantes que avaliaram a prova como curta ou muito curta em relação ao tempo total destinado à sua resolução variou de 3,6 \% na região Nordeste a 4,1 \% na região Norte.
Os percentuais, quanto ao mesmo aspecto, entre os Concluintes, foram de 4,6 \% na região Nordeste, indo até 6,0 \% na região Sudeste. 4.3 COMPREENSÃO DOS ENUNCIADOS DAS QUESTÕES 4.3.1 Componente de Formação Geral Com relação aos enunciados das questões do Componente de Formação Geral, as opiniões foram positivas, com percentuais que foram dos 78,2 \% de Concluintes aos 87,8 \% de Ingressantes que consideraram claros e objetivos os enunciados de todas ou da maioria das questões ( Gráfico 4.4 e Tabelas II.4 e II.13 no Anexo II ).
Ao se comparar as percepções dos estudantes segundo o desempenho, observou -se que, entre os Ingressantes do quarto superior, a clareza e objetividade de todos ou da maioria dos enunciados das questões foram percebidas por 90,4 \%.
Este mesmo quesito, apurado para os Concluintes, teve percentual de 79,8 \%.
No quarto de desempenho mais baixo, tal avaliação foi emitida por 82,9 \% dos Ingressantes e por 75,1 \% dos Concluintes.
A partir dessas estatísticas, é possível afirmar que todos ou a maioria dos enunciados das questões do Componente de Formação Geral estavam claros e objetivos para a maior parte dos respondentes, segundo os quartos de desempenho considerados, embora os que se situaram no quarto superior tenham declarado esta percepção com maior incidência.
Na análise regional, observa -se que as opiniões de que todos ou a maioria dos enunciados das questões do Componente de Formação Geral estavam claros e objetivos variaram entre os Ingressantes de 86,5 \% no Sul a 91,0 \%, no Centro-Oeste.
Entre os Concluintes, os valores correspondentes variaram de 76,2 \%, no Sudeste a 82,2 \%, no Centro-Oeste. 4.3.2 Componente de Conhecimento Específico Com relação aos enunciados das questões do Componente de Conhecimento Específico da prova, para os Concluintes da Área de Medicina, a clareza e a objetividade foram percebidas em todos ou na maioria destes por 86,3 \%.
Para o grupo dos Ingressantes, foram percebidas por 88,3 \%.
A proporção de estudantes que consideraram os enunciados das questões claros e objetivos foi mais elevada nos quartos superiores : 91,1 \% dos Ingressantes e 89,4 \% dos Concluintes, se comparados aos 82,7 \% e aos 82,4 \% dos mesmos grupos, respectivamente, que estiveram no quarto inferior de desempenho ( Gráfico 4.5 e Tabelas II.5 e II.14 no Anexo II ).
Nas cinco Grandes Regiões brasileiras, os percentuais de Ingressantes que consideraram claros e objetivos todos ou a maioria dos enunciados das questões do Componente de Conhecimento Específico da prova variaram de 86,9 \% ( região Norte ) a 90,3 \% ( Centro-Oeste ).
Os Concluintes tiveram índices entre 85,3 \%, na região Sudeste e 88,3 \% na região Norte. 4.4 SUFICIÊNCIA DAS INFORMAÇÕES/INSTRUÇÕES FORNECIDAS Ao avaliarem as informações/instruções fornecidas para a resolução das questões, 88,8 \% do grupo dos Ingressantes e 94,9 \% do grupo dos Concluintes respondentes de todo o Brasil afirmaram que estas eram até excessivas ou suficientes em todas ou na maioria das questões ( Gráfico 4.6 e Tabelas II.6 e II.15 no Anexo II ).
Levando -se em conta o desempenho dos participantes dos quartos inferior e superior de desempenho, como mostra o Gráfico 4.6, o percentual de participantes que avaliaram as informações/instruções como excessivas ou suficientes em todas ou na maioria das questões foi mais elevado nos quartos superiores, correspondendo a 89,6 \% dos Ingressantes e a 96,8 \% dos Concluintes, percentuais superiores à média nacional para cada um dos dois grupos.
Já nos quartos inferiores, a suficiência das informações/instruções declarada como excessiva em todas ou na maioria das questões foi percebida por 86,3 \% dos Ingressantes e por 91,8 \% dos Concluintes.
Quanto à distribuição de respondentes nas Grandes Regiões, observa -se que as proporções de estudantes que perceberam como excessivas ou suficientes em todas ou na maioria das questões as informações/instruções fornecidas para a resolução das mesmas superaram as médias nacionais, em maior medida, entre os Ingressantes da região Nordeste ( 90,1 \% ).
Entre os Concluintes, a superação da média nacional foi mais expressiva também na região Nordeste ( 95,6 \% ). 4.5 DIFICULDADE ENCONTRADA AO RESPONDER À PROVA Perguntados sobre as dificuldades com as quais se depararam ao responder à prova, 76,2 \% do grupo de Ingressantes e 11,3 \% do grupo de Concluintes apontaram o desconhecimento do conteúdo.
Os Concluintes, por sua vez, mencionaram em maior medida a forma diferente de abordagem do conteúdo como dificuldade, apontada por 31,6 \%.
Essa alternativa foi indicada por uma parcela menor dos Ingressantes, 7,8 \%.
A falta de motivação foi a dificuldade para 27,2 \% dos Concluintes e para 8,6 \% dos Ingressantes.
Considerando -se os respondentes em todo o Brasil, 5,5 \% dos Ingressantes e 25,2 \% dos Concluintes, afirmaram que não tiveram qualquer tipo de dificuldade para responder à prova ( Tabelas II.7 e II.16 no Anexo II ).
Com relação aos quartos de desempenho, o desconhecimento do conteúdo foi a dificuldade modal entre os Ingressantes, apontada por 79,1 \% daqueles no quarto superior e por 69,9 \% dos do inferior.
Entre os Concluintes, os percentuais de estudantes com a mesma percepção foi inferior ao encontrado entre os Ingressantes, 9,4 \% no quarto superior e 11,6 \% no quarto inferior.
A alternativa modal pelo grupo de Concluintes, quando desagregada pelo desempenho, foi que a dificuldade encontrada era a forma diferente de abordagem do conteúdo : 30,3 \% no quarto inferior e 32,8 \% do quarto superior assim o responderam.
Na análise por Grandes Regiões, o percentual de Ingressantes que apontaram o desconhecimento do conteúdo como dificuldade ao responder à prova variou de 72,0 \%, na região Norte a 78,9 \%, na Sudeste.
Entre os Concluintes, tal percentual foi mais baixo, variando de 10,4 \% na região Sudeste a 13,6 \%, na Centro-Oeste.
A forma diferente de abordagem do conteúdo foi citada, pelos os Concluintes, com percentuais que variaram entre 28,0 \% ( região Centro-Oeste ) e 37,1 \% ( Norte ).
Em contrapartida, os percentuais desta alternativa entre os Ingressantes foram mais baixos, variando de 6,5 \% ( região Sudeste ) a 12,4 \% ( região Norte ).
O percentual de Ingressantes que citou a falta de motivação como dificuldade superou a média nacional nas regiões Centro-Oeste ( 9,5 \% ), Sul ( 10,6 \% ) e Nordeste ( 10,1 \% ).
Entre os Concluintes, os percentuais para essa resposta variaram de 24,3 \% ( região Norte ) a 29,3 \% ( região Nordeste ).
Percentuais acima das médias nacionais dos Ingressantes das regiões Nordeste ( 6,6 \% ) e Norte ( 5,6 \% ) declararam que não tiveram qualquer dificuldade para responder à prova.
Entre os Concluintes, esta situação aconteceu nas regiões Sudeste e Sul ( ambas com 25,5 \% ), além da Centro-Oeste ( 26,2 \% ).
O Gráfico 4.7 apresenta os percentuais de estudantes, Ingressantes e Concluintes, que apontaram o desconhecimento do conteúdo como dificuldade percebida ao responder à prova. 4.6 CONTEÚDOS DAS QUESTÕES OBJETIVAS DA PROVA Ao analisarem os conteúdos das questões objetivas da prova, a maior parte dos Ingressantes, 78,5 \% do total nacional, afirmou que não estudou ainda a maioria desses conteúdos, enquanto entre os Concluintes, menos de 1,0 \% expressou a mesma opinião.
Por outro lado, a maioria dos Concluintes ( 87,4 \% ) afirmou ter estudado e aprendido muitos ou todos os conteúdos.
Para os Ingressantes, o percentual correspondente a essas alternativas foi 12,5 \% ( Gráfico 4.8 e Tabelas II.8 e II.17 no Anexo II ).
Considerando -se separadamente as opiniões de estudantes dos quartos de desempenho inferior e superior, observa -se que, para os Ingressantes, o percentual de participantes que afirmaram que não estudaram ainda a maioria desses conteúdos, foi praticamente igual para o quarto inferior e o superior : 77,7 \% e 77,2 \%, respectivamente.
Para os Concluintes do quarto inferior, 1,6 \% ofereceram essa resposta, sendo apenas 0,3 \% os estudantes do quarto superior.
Tendo em conta os quartos superiores para os dois grupos, 93,3 \% dos Concluintes e 13,8 \% dos Ingressantes, afirmaram ter estudado e aprendido muitos ou todos os conteúdos.
Na análise por Grande Região, a proporção de Ingressantes que afirmou que não estudou ainda a maioria desses conteúdos superou a média nacional nas regiões Sudeste ( 82,1 \% ), Centro-Oeste ( 79,3 \% ) e Sul ( 78,6 \% ).
Já entre os Concluintes, as regiões Sudeste ( 0,9 \% ) e Nordeste ( 1,1 \% ) apresentaram percentuais maiores do que a média nacional para esta alternativa.
Em todas as regiões, a maioria dos Concluintes afirmaram ter estudado e aprendido muitos ou todos os conteúdos, com proporções variando entre 80,7 \% na região Norte e 90,2 \% na Centro-Oeste. 4.7 TEMPO GASTO PARA CONCLUIR A PROVA Ao responderem sobre o tempo de conclusão da prova, mais da metade dos estudantes afirmou ter gasto entre duas e quatro horas.
Considerando -se separadamente as percepções do grupo de Ingressantes e do grupo de Concluintes, os percentuais, em termos nacionais, foram 72,2 \% e 78,2 \%, respectivamente ( Gráfico 4.9 e Tabelas II.9 e II.18 no Anexo II ).
Uma vez sendo analisadas as alternativas oferecidas pelos estudantes que se situaram nos quartos superior e inferior, observa -se que uma maior proporção de Ingressantes no quarto superior declarou ter gasto entre duas e quatro horas para concluir a prova quando comparados com os do quarto inferior, respectivamente 82,7 \% e 54,5 \%.
A mesma situação foi verificada no grupo dos Concluintes, no qual 84,8 \% daqueles no quarto superior e 69,8 \% daqueles no quarto inferior declaram ter gasto entre duas e quatro horas para concluir a prova.
Nas cinco Grandes Regiões brasileiras, os Ingressantes que utilizaram entre duas e quatro horas para finalizar a prova superaram o percentual nacional nas regiões Norte ( 74,7 \% ), Nordeste ( 75,7 \% ) e Centro-Oeste ( 76,9 \% ).
Entre os Concluintes, o mesmo foi verificado entre os estudantes das regiões Sul ( 79,5 \% ) e Centro-Oeste ( 79,9 \% ), como mostra o Gráfico 4.9.
A região Sul foi a que apresentou o menor percentual de Ingressantes que dispensaram entre duas e quatro horas para concluir a prova : 68,7 \%.
Entre os Concluintes, foi a região Sudeste com 77,8 \%, que esteve na mesma situação.

% matched lemmas: 
\end{textsample}
